\chapter{Blocks 1--2: Unix Foundations}\label{chap:unix}

\textbf{Primary:} Command line, text processing, bash scripting\quad|\quad\textbf{Interleave:} Git, project setup

\begin{brainbox}[Why this matters]
Every bioinformatics pipeline runs on Unix. Genomic data is text. If you can wrangle text in a terminal, you can wrangle genomes.
\end{brainbox}

\begin{tipbox}[Need to install something?]
Every non-default tool named anywhere in this curriculum has install commands (macOS and Linux) and an official link in \textbf{\hyperref[app:tools]{Appendix~A}}. Flip to it the moment you hit a command that isn't on your machine --- don't guess.
\end{tipbox}


\section{Block 1: The Command Line}

\sprint{1}{Navigate \& Look Around}{25}
\begin{itemize}[nosep]
  \item Open your terminal. Run: \texttt{pwd}, \texttt{ls}, \texttt{ls -la}, \texttt{cd \textasciitilde}, \texttt{cd /}, \texttt{cd -}, \texttt{tree} (install with conda if needed)
  \item Make a project directory: \texttt{mkdir -p \textasciitilde/bioinfo/block1/\{data,scripts,results\}}
  \item Explore: what's in \texttt{/usr/bin}? How many files? (\texttt{ls /usr/bin | wc -l})
\end{itemize}
\begin{winbox}[Checkpoint]
Screenshot your project directory tree. Paste it in journal.md.
\end{winbox}

\sprint{2}{Pipes --- The Unix Superpower}{25}
\begin{itemize}[nosep]
  \item The pipe (\texttt{|}) sends output from one command into another. This is the entire Unix philosophy.
  \item Try: \texttt{echo 'ATGCGATCGATCG' | grep -o . | sort | uniq -c | sort -rn} --- count nucleotide frequencies in one line
  \item Download a small FASTA file:
    \begin{lstlisting}
curl -O https://downloads.yeastgenome.org/sequence/
  S288C_reference/genome_releases/
  S288C_reference_genome_Current_Release.tgz
tar xzf S288C_reference_genome_Current_Release.tgz
    \end{lstlisting}
  \item Explore it: \texttt{grep '>' *.fsa | head} --- what are these header lines telling you?
\end{itemize}
\begin{winbox}[Checkpoint]
You piped three commands together and saw structured output.
\end{winbox}

\sprint{3}{grep + awk = Data Surgery}{25}
\begin{itemize}[nosep]
  \item \texttt{grep}: searches for patterns. \texttt{grep -c '>' file.fasta} counts sequences. \texttt{grep -v} inverts (excludes matches).
  \item \texttt{awk}: extracts and transforms columns. \texttt{awk -F'\textbackslash t' '\{print \$1, \$2, \$5\}'} pulls chr, pos, and qual from a VCF.
  \item Grab a small VCF to practice on. Yeast works well --- a whole-organism variant file that's only a few MB:
    \begin{lstlisting}
curl -O https://ftp.ensembl.org/pub/
  current_variation/vcf/saccharomyces_cerevisiae/
  saccharomyces_cerevisiae.vcf.gz
gunzip saccharomyces_cerevisiae.vcf.gz
    \end{lstlisting}
    If that link has moved, browse \texttt{https://ftp.ensembl.org/pub/current\_variation/vcf/} and pick any small organism's VCF. Other good ``small VCF'' sources: the bcftools test data (\texttt{github.com/samtools/bcftools}, \texttt{test/} directory) or a single-chromosome 1000 Genomes file restricted to a small region with \texttt{tabix}.
  \item Count variants per chromosome with:
    \begin{lstlisting}
grep -v '^#' file.vcf | awk '{print $1}' | \
  sort | uniq -c | sort -rn
    \end{lstlisting}
\end{itemize}
\begin{winbox}[Checkpoint]
You summarized a genome's worth of variants in one line.
\end{winbox}

\sprint{4}{Put It All Together --- Pick a Restriction Enzyme for Cloning}{25}
A real scenario: you want to shotgun-clone $\sim$3\,kb fragments of the yeast genome into a plasmid. Three cheap 6-cutters are on the shelf --- EcoRI (\texttt{GAATTC}), BamHI (\texttt{GGATCC}), and HindIII (\texttt{AAGCTT}). \textit{Before} ordering anything, predict from the genome sequence how often each enzyme cuts and which one produces fragments closest to your 3\,kb target. Do it all with pipes --- no scripts yet. Paste every command and its output into \texttt{journal.md} as you go.
\begin{enumerate}[nosep]
  \item \textbf{Confirm and organize.} \texttt{head} the yeast FASTA --- the first line should start with \texttt{>}. Then move the FASTA (and the VCF from Sprint 3) into \texttt{data/} in one pipeline: \texttt{ls *.fsa *.vcf 2>/dev/null | xargs -I\{\} mv \{\} \textasciitilde/bioinfo/block1/data/}. Verify with \texttt{ls \textasciitilde/bioinfo/block1/data/}.
  \item \textbf{Measure the genome.} Strip headers and newlines, count characters:
    \begin{lstlisting}
grep -v '>' ~/bioinfo/block1/data/*.fsa | \
  tr -d '\n' | wc -c
    \end{lstlisting}
    For S288C, expect roughly 12\,Mb. Save this number --- you'll divide by it next.
  \item \textbf{Count cut sites per enzyme.} Template for EcoRI:
    \begin{lstlisting}
grep -v '>' ~/bioinfo/block1/data/*.fsa | \
  tr -d '\n' | grep -o GAATTC | wc -l
    \end{lstlisting}
    Run it three times --- once each for \texttt{GAATTC}, \texttt{GGATCC}, \texttt{AAGCTT}. Which cuts most? Fewest?
  \item \textbf{Why one strand is enough.} All three recognition sites are palindromes (read the same on both strands), so every real cut site is already captured by a forward-strand match. For non-palindromic motifs you'd also have to scan the reverse complement --- good to know before you generalize this approach.
  \item \textbf{Compute average fragment size.} Fragment size $\approx$ genome length $\div$ number of cut sites. Use \texttt{bc}: \texttt{echo "12157105 / NSITES" | bc} (swap in your actual genome length and cut count). Expect numbers in the range of hundreds of bp to tens of kb.
  \item \textbf{Pick your enzyme.} Which enzyme's average fragment size is closest to your 3\,kb target? Write one sentence in \texttt{journal.md} explaining your choice, and flag which enzyme would be a bad idea (cuts so often you'd get tiny shredded fragments, or so rarely the fragments are too big to clone).
\end{enumerate}
\begin{winbox}[Checkpoint]
Your journal shows the \texttt{head} check, the move-into-\texttt{data/} pipeline, the genome length, cut-site counts for three enzymes, their average fragment sizes, and your enzyme pick with a one-sentence justification --- all produced with pipes.
\end{winbox}

\begin{tipbox}[Pattern to remember]
In Unix bioinformatics, the workflow is almost always: EXTRACT columns (cut/awk) $\to$ FILTER rows (grep) $\to$ SORT $\to$ COUNT (uniq -c). This pattern covers 80\% of quick data exploration.
\end{tipbox}


\section{Block 2: Scripts, Environments \& Git}

\sprint{5}{Your First Bash Script}{25}
\begin{itemize}[nosep]
  \item Create a file: \texttt{nano \textasciitilde/bioinfo/scripts/count\_variants.sh}
  \item Add: \texttt{\#!/usr/bin/env bash}, then \texttt{set -euo pipefail} (this catches errors immediately)
  \item Write a script that takes a VCF filename as \texttt{\$1} and outputs a per-chromosome variant count table
  \item Make it executable: \texttt{chmod +x count\_variants.sh}, then run it
\end{itemize}
\begin{winbox}[Checkpoint]
You have a reusable script. Run it on a different VCF. Same script, different data.
\end{winbox}

\sprint{6}{Conda \& Mamba: Isolated, Reproducible Environments}{25}
\begin{itemize}[nosep]
  \item Create an environment: \texttt{conda create -n biotools samtools bcftools bedtools fastp}
  \item Activate it: \texttt{conda activate biotools}. Verify: \texttt{samtools --version}.
  \item Export it: \texttt{conda env export > environment.yml} --- anyone can now recreate your setup with \texttt{conda env create -f environment.yml}.
  \item Deactivate: \texttt{conda deactivate}. Run \texttt{samtools} again --- it's no longer available. Environments isolate dependencies from each other and from the system.
  \item \textbf{Use mamba for speed.} \texttt{conda}'s dependency solver is slow on large environments. \texttt{mamba} is a drop-in replacement (C++ solver) that resolves the same specs in seconds. Miniforge (\url{https://github.com/conda-forge/miniforge}) ships with it; otherwise \texttt{conda install -n base -c conda-forge mamba}. Swap it in wherever you'd type \texttt{conda}: \texttt{mamba create -n biotools samtools bcftools bedtools fastp}.
\end{itemize}
\begin{warnbox}[Cross-architecture caveats]
Conda packages are built per (OS, CPU) pair --- \texttt{linux-64}, \texttt{osx-64}, \texttt{osx-arm64} (Apple Silicon), \texttt{linux-aarch64} (ARM Linux / some HPCs). Not every bioinformatics package has a build for every platform. On an M1/M2/M3 Mac, a missing \texttt{osx-arm64} build usually surfaces as \texttt{PackagesNotFoundError}; common fixes are to (1) create the env under Rosetta with \texttt{CONDA\_SUBDIR=osx-64 conda create \ldots}, or (2) find an alternative package with native ARM support. An \texttt{environment.yml} built on one platform is not guaranteed to solve on another --- expect small fix-ups when moving between laptop and cluster, and note the architecture in a comment at the top of the file.
\end{warnbox}
\begin{winbox}[Checkpoint]
Your \texttt{environment.yml} is a reproducible recipe. Commit it to Git, with a header comment noting which architecture built it.
\end{winbox}

\sprint{7}{Containers --- Reproducibility Beyond Conda}{25}
\begin{brainbox}[What is a container?]
A container is a single image file that bundles an entire Linux userland --- OS libraries, the tool you care about, and everything the tool depends on --- so it runs bit-for-bit identically on any compatible host. Conda manages packages \textit{inside} ecosystems (Python, R, bioconda); containers freeze the layer \textit{below} that, where system libraries, \texttt{glibc} versions, and compiled binaries live. A container is not a virtual machine --- there's no guest kernel --- so startup is near-instant and overhead is minimal.

Rule of thumb: use conda for interactive development (fast to edit, easy to iterate). Use a container when you ship a pipeline to a cluster, share an analysis with a collaborator on a different OS, or archive a workflow alongside a manuscript --- situations where ``works on my machine'' has to become ``works on yours.''
\end{brainbox}

\begin{itemize}[nosep]
  \item \textbf{Two runtimes you'll encounter.} \texttt{Docker} is the developer-facing tool --- it requires a privileged daemon, so most shared HPCs will not run it. Install Docker Desktop on macOS (\url{https://docs.docker.com/desktop/install/mac-install/}) or the Engine on Linux (\url{https://docs.docker.com/engine/install/}). \texttt{Apptainer} (formerly \texttt{Singularity}) is the HPC-safe equivalent: no daemon, no root, and a container is a single \texttt{.sif} file you can \texttt{scp} around. HPCs usually provide a module (\texttt{module avail apptainer}); for a personal Linux install, see \url{https://apptainer.org/docs/user/main/quick_start.html}.
  \item \textbf{Pull a biocontainer.} The Biocontainers project auto-builds one container per bioconda package. Try samtools:
    \begin{lstlisting}
docker pull \
  quay.io/biocontainers/samtools:1.19--h50ea8bc_0
    \end{lstlisting}
  \item \textbf{Run without installing.} Execute the container's \texttt{samtools} with no other setup:
    \begin{lstlisting}
docker run --rm \
  quay.io/biocontainers/samtools:1.19--h50ea8bc_0 \
  samtools --version
    \end{lstlisting}
    \texttt{--rm} discards the stopped container after it exits; the image stays cached for next time.
  \item \textbf{Mount your data.} Containers see an isolated filesystem by default. Map your current directory in so the tool can read local files:
    \begin{lstlisting}
docker run --rm -v "$PWD":/work -w /work \
  quay.io/biocontainers/samtools:1.19--h50ea8bc_0 \
  samtools view file.bam | head
    \end{lstlisting}
    \texttt{-v HOST:CONTAINER} mounts a host directory; \texttt{-w} sets the working directory inside.
  \item \textbf{On an HPC, use Apptainer.} Same image, no daemon required:
    \begin{lstlisting}
apptainer pull \
  docker://quay.io/biocontainers/samtools:1.19--h50ea8bc_0
apptainer exec samtools_1.19--h50ea8bc_0.sif \
  samtools --version
    \end{lstlisting}
    Apptainer auto-mounts your home directory and the current working directory, so most pipelines just work.
\end{itemize}
\begin{tipbox}[Conda vs.\ container --- when to reach for which]
Conda gives you speed and easy iteration at the cost of reproducibility guarantees --- two users with the same \texttt{environment.yml} may still resolve to different package builds on different platforms. A container gives you a fixed binary that will run identically for years, at the cost of being heavier to build and harder to edit. Many published pipelines use both: conda during development, a container (often built \textit{from} the conda env) for the frozen release.
\end{tipbox}
\begin{winbox}[Checkpoint]
You pulled a container, ran a bioinformatics tool from it without installing anything locally, and mounted a host directory into it. You know which runtime (Docker vs.\ Apptainer) you'd reach for on a laptop vs.\ an HPC, and why.
\end{winbox}

\sprint{8}{Git + GitHub --- Your Time Machine and Portfolio}{25}
\begin{itemize}[nosep]
  \item \textbf{Initialize and commit.}
    \begin{lstlisting}
cd ~/bioinfo && git init -b main && \
  git add -A && \
  git commit -m 'Blocks 1-2: Unix foundations'
    \end{lstlisting}
  \item Edit something. \texttt{git diff} to inspect, \texttt{git add}/\texttt{git commit} to save, \texttt{git log --oneline} to review.
  \item \textbf{.gitignore.} Exclude anything big or generated: \texttt{*.bam}, \texttt{*.fastq.gz}, \texttt{*.vcf.gz}, \texttt{*.sam}. Never commit raw data --- commit the code that produces it and a pointer to where it lives.
  \item \textbf{Set up an SSH key for GitHub} (one-time; after this, no password prompts):
    \begin{lstlisting}
ssh-keygen -t ed25519 -C "you@example.com"
# press Enter to accept the default path;
# optionally set a passphrase
cat ~/.ssh/id_ed25519.pub
    \end{lstlisting}
    Copy the printed public key, then go to github.com $\to$ Settings $\to$ SSH and GPG keys $\to$ New SSH key and paste it. Test: \texttt{ssh -T git@github.com} should greet you by username.
  \item \textbf{Publish your repo.} Create an empty repo on github.com (don't add a README --- your local copy already has files). Then from your local repo:
    \begin{lstlisting}
git remote add origin \
  git@github.com:<you>/bioinfo-journey.git
git push -u origin main
    \end{lstlisting}
    The \texttt{-u} flag wires \texttt{main} to the remote so future \texttt{git push}es need no arguments.
  \item \textbf{Clone and pull.} To download a repo onto a new machine or grab someone else's work: \texttt{git clone git@github.com:user/repo.git}. Inside an existing clone, \texttt{git pull} fetches new commits from the remote and merges them into your branch --- this is how you keep a checkout in sync. Practice: clone your own repo into \texttt{/tmp/}, push a change from your main checkout, then run \texttt{git pull} in \texttt{/tmp/} and confirm the change arrived.
\end{itemize}
\begin{winbox}[Checkpoint]
Your bioinfo-journey repo has real commits, ignores big files, pushes over SSH without password prompts, and you've proved you can \texttt{clone} and \texttt{pull} from it. This is your portfolio.
\end{winbox}

\sprint{9}{Integration Sprint: Full Pipeline}{50}
\begin{itemize}[nosep]
  \item Combine everything: write a bash script that creates a conda env, downloads a FASTQ from SRA (use \texttt{fastq-dump} or \texttt{fasterq-dump} from the SRA Toolkit --- \texttt{mamba install -c bioconda sra-tools}), runs \texttt{fastp} for QC, and logs results to a file.
  \item Use \texttt{tmux} (macOS: \texttt{brew install tmux}; Linux: \texttt{sudo apt install tmux}) or \texttt{screen} so you can disconnect and reconnect (crucial for HPC).
  \item Commit the script and the QC results (not the raw data) to Git.
\end{itemize}
\begin{winbox}[Checkpoint]
A complete, reproducible, version-controlled QC pipeline. From nothing to this in two blocks.
\end{winbox}

\begin{restbox}[End of Blocks 1--2]
Take a day to step away before starting Python. The command-line fundamentals are now in place.
\end{restbox}

\subsection{Can I move on? Checkpoint:}
\begin{itemize}[nosep]
  \item You can navigate any directory, find files, and chain 3+ commands with pipes
  \item You can write a bash script that processes a genomic file and logs output
  \item You have a conda environment with bioinformatics tools installed
  \item Your work is in a Git repo on GitHub
\end{itemize}
\textit{If not, spend 2--3 more days on the gaps. No rush.}

\subsection{Resources}
\begin{itemize}[nosep]
  \item Software Carpentry: The Unix Shell (free, interactive)
  \item Buffalo -- \textit{Bioinformatics Data Skills}, Ch.\ 3, 7, 12
  \item explainshell.com --- paste any command and see what each part does
\end{itemize}
